% q0027.tex — Follow the leader (Stackelberg vs Cournot) [Numeric Answer]
\question{
  id        = {0027},
  title     = {Follow the Leader},
  source    = {OBECON},
  year      = {2026},
  round     = {Seletiva IEO -- Phase A},
  language  = {English},
  topic     = {Game Theory},
  subtopic  = {Industrial Organization / Stackelberg vs Cournot},
  type      = {numeric},
  statement = {%
    Two oil companies, $A$ and $B$, initially compete in a duopoly under
    Cournot competition. To reduce the negative externalities of oil
    production, the government introduces a new policy that changes the market
    structure: in each period, firms compete following the Stackelberg model,
    with one firm acting as the leader and the other as the follower. Firm~$A$
    will be the leader in the first period. Then, leadership rotates according
    to the following rule: the firm producing the \emph{smaller quantity} in a
    given period becomes the leader in the next period (if they produce the
    same quantity, the previous follower becomes the next leader). Production
    quantities must be non-negative integers.

    The market demand is given by $Q = 360 - p$. Firms $A$ and $B$ both have
    constant marginal costs of 120, with no fixed costs. The interest rate is
    20\%. Both firms are rational, the game extends over an infinite horizon,
    and firms never cooperate or collude (no strategies contingent on past
    actions).

    The government wishes to assess whether this policy effectively reduces
    total production. Considering the total output produced over the next 10
    years under this new policy, what is the difference compared to the output
    that would arise over the same horizon if the firms remained in Cournot
    competition?

    Give the value of $(S_A + S_B) - (C_A + C_B)$, where $S_i$ is the output
    of firm~$i$ in the next 10 years under the new rule, and $C_i$ is the
    output of firm~$i$ under Cournot competition.
  },
  choices   = {},
  answer    = {},
  solution  = {%
    % TODO: add solution
  },
}
